\section{A key lemma}
\label{section-key-lemma}

\noindent
In this section we apply the results above to prove
Lemma \ref{lemma-milnor-gersten-low-degree}.
This lemma is a low degree case of the statement
that there is a complex for Milnor K-theory similar
to the Gersten-Quillen complex in Quillen's K-theory.
See Remark \ref{remark-gersten-complex-milnor}.

\begin{lemma}
\label{lemma-perpare-key}
Let $(A, \mathfrak m)$ be a $2$-dimensional Noetherian local ring.
Let $t \in \mathfrak m$ be a nonzerodivisor. Say
$V(t) = \{\mathfrak m, \mathfrak q_1, \ldots, \mathfrak q_r\}$.
Let $A_{\mathfrak q_i} \subset B_i$ be a finite ring
extension with $B_i/A_{\mathfrak q_i}$ annihilated by a power of
$t$. Then there exists a finite extension $A \subset B$ of
local rings identifying residue fields
with $B_i \cong B_{\mathfrak q_i}$ and $B/A$ annihilated
by a power of $t$.
\end{lemma}

\begin{proof}
Choose $n > 0$ such that $B_i \subset t^{-n}A_{\mathfrak q_i}$.
Let $M \subset t^{-n}A$, resp.\ $M' \subset t^{-2n}A$ be the
$A$-submodule consisting of elements mapping to $B_i$ in
$t^{-n}A_{\mathfrak q_i}$, resp.\ $t^{-2n}A_{\mathfrak q_i}$.
Then $M \subset M'$ are finite $A$-modules as $A$ is Noetherian
and $M_{\mathfrak q_i} = M'_{\mathfrak q_i} = B_i$ as localization
is exact. Thus $M'/M$ is annihilated by $\mathfrak m^c$ for some
$c > 0$. Observe that $M \cdot M \subset M'$ under the multiplication
$t^{-n}A \times t^{-n}A \to t^{-2n}A$. Hence
$B = A + \mathfrak m^{c + 1}M$ is a finite $A$-algebra with the correct
localizations. We omit the verification that $B$ is local with
maximal ideal $\mathfrak m + \mathfrak m^{c + 1}M$.
\end{proof}

\begin{lemma}
\label{lemma-key-nonzerodivisors}
Let $(A, \mathfrak m)$ be a $2$-dimensional Noetherian local ring.
Let $a, b \in A$ be nonzerodivisors.
Then we have
$$
\sum
\text{ord}_{A/\mathfrak q}(\partial_{A_{\mathfrak q}}(a, b))
=
0
$$
where the sum is over the height $1$ primes $\mathfrak q$ of $A$.
\end{lemma}

\begin{proof}
If $\mathfrak q$ is a height $1$ prime of $A$ such that $a, b$
map to a unit of $A_\mathfrak q$, then $\partial_{A_\mathfrak q}(a, b) = 1$.
Thus the sum is finite. In fact, if
$V(ab) = \{\mathfrak m, \mathfrak q_1, \ldots, \mathfrak q_r\}$
then the sum is over $i = 1, \ldots, r$.
For each $i$ we pick an extension $A_{\mathfrak q_i} \subset B_i$
as in Lemma \ref{lemma-prepare-tame-symbol} for $a, b$.
By Lemma \ref{lemma-perpare-key} with $t = ab$ and the given list of primes
we may assume we have a finite local extension $A \subset B$
with $B/A$ annihilated by a power of $ab$ and such that
for each $i$ the $B_{\mathfrak q_i} \cong B_i$.
Observe that if $\mathfrak q_{i, j}$ are the primes of $B$
lying over $\mathfrak q_i$ then we have
$$
\text{ord}_{A/\mathfrak q_i}(\partial_{A_{\mathfrak q_i}}(a, b))
=
\sum\nolimits_j
\text{ord}_{B/\mathfrak q_{i, j}}(\partial_{B_{\mathfrak q_{i, j}}}(a, b))
$$
by Lemma \ref{lemma-norm-down-tame-symbol} and
Algebra, Lemma \ref{algebra-lemma-finite-extension-dim-1}.
Thus we may replace $A$ by $B$ and
reduce to the case discussed in the next paragraph.

\medskip\noindent
Assume for each $i$ there is a nonzerodivisor
$\pi_i \in A_{\mathfrak q_i}$ and units $u_i, v_i \in A_{\mathfrak q_i}$
such that for some integers $e_i, f_i \geq 0$ we have
$$
a = u_i \pi_i^{e_i},\quad b = v_i \pi_i^{f_i}
$$
in $A_{\mathfrak q_i}$. Setting
$m_i = \text{length}_{A_{\mathfrak q_i}}(A_{\mathfrak q_i}/\pi_i)$
we have
$\partial_{A_{\mathfrak q_i}}(a, b) =
((-1)^{e_if_i}u_i^{f_i}v_i^{-e_i})^{m_i}$ by definition.
Since $a, b$ are nonzerodivisors the
$(2, 1)$-periodic complex $(A/(ab), a, b)$ has vanishing cohomology.
Denote $M_i$ the image of $A/(ab)$ in $A_{\mathfrak q_i}/(ab)$.
Then we have a map
$$
A/(ab) \longrightarrow \bigoplus M_i
$$
whose kernel and cokernel are supported in $\{\mathfrak m\}$
and hence have finite length. Thus we see that
$$
\sum e_A(M_i, a, b) = 0
$$
by Lemma \ref{lemma-compare-periodic-lengths}. Hence it suffices to show
$e_A(M_i, a, b) =
- \text{ord}_{A/\mathfrak q_i}(\partial_{A_{\mathfrak q_i}}(a, b))$.

\medskip\noindent
Let us prove this first, in case $\pi_i, u_i, v_i$ are the images
of elements $\pi_i, u_i, v_i \in A$ (using the same symbols
should not cause any confusion). In this case we get
\begin{align*}
e_A(M_i, a, b) & =
e_A(M_i, u_i\pi_i^{e_i}, v_i\pi_i^{f_i}) \\
& =
e_A(M_i, \pi_i^{e_i}, \pi_i^{f_i}) -
e_A(\pi_i^{e_i}M_i, 0, u_i) +
e_A(\pi_i^{f_i}M_i, 0, v_i) \\
& =
0 -
f_im_i\text{ord}_{A/\mathfrak q_i}(u_i) +
e_im_i\text{ord}_{A/\mathfrak q_i}(v_i) \\
& =
-m_i\text{ord}_{A/\mathfrak q_i}(u_i^{f_i}v_i^{-e_i}) =
-\text{ord}_{A/\mathfrak q_i}(\partial_{A_{\mathfrak q_i}}(a, b))
\end{align*}
The second equality holds by Lemma \ref{lemma-multiply-period-length}.
Observe that
$M_i \subset (M_i)_{\mathfrak q_i} = A_{\mathfrak q_i}/(\pi_i^{e_i + f_i})$
and
$(\pi_i^{e_i}M_i)_{\mathfrak q_i} \cong A_{\mathfrak q_i}/\pi_i^{f_i}$ and
$(\pi_i^{f_i}M_i)_{\mathfrak q_i} \cong A_{\mathfrak q_i}/\pi_i^{e_i}$.
The $0$ in the third equality comes from
Lemma \ref{lemma-powers-period-length-zero}
and the other two terms come from
Lemma \ref{lemma-length-multiplication}.
The last two equalities follow from multiplicativity of
the order function and from the definition of our tame symbol.

\medskip\noindent
In general, we may first choose $c \in A$, $c \not \in \mathfrak q_i$
such that $c\pi_i \in A$. After replacing $\pi_i$ by $c\pi_i$
and $u_i$ by $c^{-e_i}u_i$ and $v_i$ by $c^{-f_i}v_i$
we may and do assume $\pi_i$ is in $A$.
Next, choose an $c \in A$, $c \not \in \mathfrak q_i$
with $cu_i, cv_i \in A$. Then we observe that
$$
e_A(M_i, ca, cb) = e_A(M_i, a, b) - e_A(aM_i, 0, c) + e_A(bM_i, 0, c)
$$
by Lemma \ref{lemma-length-multiplication}.
On the other hand, we have
$$
\partial_{A_{\mathfrak q_i}}(ca, cb) =
c^{m_i(f_i - e_i)}\partial_{A_{\mathfrak q_i}}(a, b)
$$
in $\kappa(\mathfrak q_i)^*$ because $c$ is a unit in $A_{\mathfrak q_i}$.
The arguments in the previous paragraph show that
$e_A(M_i, ca, cb) = -
\text{ord}_{A/\mathfrak q_i}(\partial_{A_{\mathfrak q_i}}(ca, cb))$.
Thus it suffices to prove
$$
e_A(aM_i, 0, c) = \text{ord}_{A/\mathfrak q_i}(c^{m_if_i})
\quad\text{and}\quad
e_A(bM_i, 0, c) = \text{ord}_{A/\mathfrak q_i}(c^{m_ie_i})
$$
and this follows from Lemma \ref{lemma-length-multiplication}
by the description (see above)
of what happens when we localize at $\mathfrak q_i$.
\end{proof}

\begin{lemma}[Key Lemma]
\label{lemma-milnor-gersten-low-degree}
\begin{reference}
When $A$ is an excellent ring this is \cite[Proposition 1]{Kato-Milnor-K}.
\end{reference}
Let $A$ be a $2$-dimensional Noetherian local domain with fraction field $K$.
Let $f, g \in K^*$.
Let $\mathfrak q_1, \ldots, \mathfrak q_t$ be the height
$1$ primes $\mathfrak q$ of $A$ such that either $f$ or $g$ is not an
element of $A^*_{\mathfrak q}$.
Then we have
$$
\sum\nolimits_{i = 1, \ldots, t}
\text{ord}_{A/\mathfrak q_i}(\partial_{A_{\mathfrak q_i}}(f, g))
=
0
$$
We can also write this as
$$
\sum\nolimits_{\text{height}(\mathfrak q) = 1}
\text{ord}_{A/\mathfrak q}(\partial_{A_{\mathfrak q}}(f, g))
=
0
$$
since at any height $1$ prime $\mathfrak q$
of $A$ where $f, g \in A^*_{\mathfrak q}$
we have $\partial_{A_{\mathfrak q}}(f, g) = 1$.
\end{lemma}

\begin{proof}
Since the tame symbols $\partial_{A_{\mathfrak q}}(f, g)$ are
bilinear and the order functions $\text{ord}_{A/\mathfrak q}$
are additive it suffices to prove the formula when
$f$ and $g$ are elements of $A$. This case is proven in
Lemma \ref{lemma-key-nonzerodivisors}.
\end{proof}

\begin{remark}[Milnor K-theory]
\label{remark-gersten-complex-milnor}
For a field $k$ let us denote $K^M_*(k)$ the quotient of
the tensor algebra on $k^*$ divided by the two-sided ideal
generated by the elements $x \otimes 1 - x$ for $x \in k \setminus \{0, 1\}$.
Thus $K^M_0(k) = \mathbf{Z}$, $K_1^M(k) = k^*$, and
$$
K^M_2(k) = k^* \otimes_\mathbf{Z} k^* / \langle x \otimes 1 - x \rangle
$$
If $A$ is a discrete valuation ring with fraction field $F = \text{Frac}(A)$
and residue field $\kappa$, there is a tame symbol
$$
\partial_A : K_{i + 1}^M(F) \to K_i^M(\kappa)
$$
defined as in Section \ref{section-tame-symbol}; see \cite{Kato-Milnor-K}.
More generally, this map can be extended to the case where $A$ is an
excellent local domain of dimension $1$ using normalization and norm
maps on $K_i^M$, see \cite{Kato-Milnor-K}; presumably the method in
Section \ref{section-tame-symbol} can be used to extend the construction
of the tame symbol $\partial_A$ to arbitrary Noetherian local domains $A$
of dimension $1$. Next, let $X$ be a Noetherian scheme with a
dimension function $\delta$. Then we can use these tame symbols to get
the arrows in the following:
$$
\bigoplus\nolimits_{\delta(x) = j + 1} K^M_{i + 1}(\kappa(x))
\longrightarrow
\bigoplus\nolimits_{\delta(x) = j} K^M_i(\kappa(x))
\longrightarrow
\bigoplus\nolimits_{\delta(x) = j - 1} K^M_{i - 1}(\kappa(x))
$$
However, it is not clear, that the composition is zero, i.e., that
we obtain a complex of abelian groups.
For excellent $X$ this is shown in \cite{Kato-Milnor-K}.
When $i = 1$ and $j$ arbitrary, this follows from
Lemma \ref{lemma-milnor-gersten-low-degree}.
\end{remark}





















